% ============================================================================
%  Tehnički vodič kroz kod — Detekcija klikbejt naslova (OPJ 2025/2026)
%  Kompajliranje:  tectonic Tehnicki-vodic-kroz-kod.tex
% ============================================================================
\documentclass[11pt,a4paper]{article}

\usepackage{fontspec}            % XeTeX/tectonic: nativni Unicode (radi i u Verbatim)
\setmonofont{lmmono10-regular.otf}[BoldFont=lmmonolt10-bold.otf]  % monospace, Unicode
\usepackage[serbian]{babel}
\usepackage[margin=2.3cm]{geometry}
\usepackage{amsmath}             % \tfrac, \cdot...
\usepackage{xcolor}
\usepackage{fancyvrb}
\usepackage{booktabs}
\usepackage{array}
\usepackage{enumitem}
\usepackage[most]{tcolorbox}
\usepackage{hyperref}
\hypersetup{colorlinks=true, linkcolor=black, urlcolor=blue, citecolor=black}

\definecolor{codebg}{RGB}{245,245,245}
\definecolor{coderule}{RGB}{200,200,200}

% „Teorija" okvir — objašnjenje pojma uz kod
\newtcolorbox{teorija}[1][]{
  colback=blue!4, colframe=blue!45!black, fonttitle=\bfseries\small,
  title={\faTheory Teorija#1}, boxrule=0.4pt, breakable,
  left=5pt, right=5pt, top=3pt, bottom=3pt, fontupper=\small}
\newcommand{\faTheory}{}  % bez ikone (bez dodatnih font-paketa)

% okvir za kod (verbatim, podržava UTF-8 dijakritike)
\DefineVerbatimEnvironment{code}{Verbatim}%
  {frame=single, rulecolor=\color{coderule}, fontsize=\small,
   xleftmargin=6pt, xrightmargin=6pt, framesep=5pt, baselinestretch=0.95}

\newcommand{\fajl}[1]{\textbf{\texttt{#1}}}
\newcommand{\kod}[1]{\texttt{#1}}

\title{\textbf{Tehnički vodič kroz kod}\\[2pt]
\large Detekcija klikbejt naslova u sportskim vestima na srpskom jeziku}
\author{Obrada prirodnih jezika 2025/2026 --- Elektrotehnički fakultet}
\date{}

\begin{document}
\maketitle

\noindent
Ovaj dokument prolazi kroz \emph{ceo izvorni kod} projekta, fajl po fajl, i za svaki
objašnjava \textbf{šta radi}, \textbf{kako je tehnički realizovan} i \textbf{šta
proizvodi}, uz primere gde pomažu razumevanju. Namenjen je za pripremu odbrane --- da
se kroz projekat može proći uporedo sa kodom. Pregled „gde se šta nalazi'' (kraća
verzija) je u \texttt{docs/Mapa-koda.md}, a komande za pokretanje u
\texttt{docs/Komande.md}.

\tableofcontents
\vspace{4mm}\hrule

% ===========================================================================
\section{Arhitektura i tok podataka}
Projekat je organizovan oko \textbf{četiri faze}, gde izlaz jedne faze ulazi u
sledeću. Sav kod je u Python-u, podeljen u pakete pod \kod{src/} po fazama.

\begin{code}
Faza 1  scraping        ==>  data/raw/*.csv            (sirovi naslovi)
Faza 2  anotacija       ==>  data/annotated/dataset.tsv (1100/1100) + IAA (kappa)
Faza 3  modeli          ==>  results/*.csv             (baseline / enkoder / dekoder)
Faza 4  izveštaj        ==>  report/tables/*.tex + figures/*.png ==> izvestaj.pdf
\end{code}

\noindent
Ključna ideja koja se ponavlja kroz kod: \emph{determinizam i reproduktivnost}
(fiksiran \kod{seed=42} svuda gde se uzorkuje/meša) i \emph{razdvojene zavisnosti po
fazi} (\kod{requirements/requirements-*.txt}), da bi se svaka faza mogla pokrenuti
nezavisno.

% ===========================================================================
\section{Faza 1 --- Prikupljanje podataka (\texttt{src/scraping/})}

\subsection{\texttt{sportklub\_api.py} --- glavni skraper}
Prikuplja naslove sa portala SportKlub preko njegovog \textbf{WordPress REST API-ja}.
Umesto „grebanja'' HTML stranica, koristi se zvanični JSON endpoint
\kod{/wp-json/wp/v2/posts}, koji vraća čist naslov, datum i kategorije --- mnogo
pouzdanije od parsiranja HTML-a.

\textbf{Kako radi (paginacija):} u petlji se traži stranica po stranica (po 100
naslova), dok se ne sakupi traženi broj. Između zahteva stoji pauza (\kod{time.sleep})
da se API ne preoptereti:

\begin{code}
API = "https://sportklub.n1info.rs/wp-json/wp/v2"
while len(rows) < broj:
    params = {"per_page": 100, "page": page,
              "_fields": "title,link,date,categories"}
    r = requests.get(f"{API}/posts", headers=HEADERS, params=params, timeout=20)
    ...
    page += 1
    time.sleep(pause)        # ljubazno prema serveru
\end{code}

\textbf{Čišćenje naslova.} WordPress vraća naslove sa HTML entitetima (npr.
\kod{\&\#8220;} umesto navodnika). Funkcija \kod{\_clean\_html} ih dekodira i uklanja
zaostale tagove:

\begin{code}
def _clean_html(s):
    s = html.unescape(s)          # &#8220; -> „
    s = re.sub(r"<[^>]+>", "", s)  # ukloni <...> tagove
    return s.strip()
\end{code}

\textbf{Metapodaci.} Svaki naslov se pakuje u \kod{Headline} strukturu sa poljima:
naslov, izvor, URL, datum objave, rubrika (ime kategorije, dobijeno iz
\kod{fetch\_categories}) i datum prikupljanja --- čime je ispunjen zahtev da se
zabeleži jedinstveni izvor/URL.

\textbf{Izlaz:} \kod{data/raw/sportklub\_api\_<datum>.csv}. Namerno se povlači
\emph{višak} (3000--4000) da bi se posle anotacije skup mogao izbalansirati.

\subsection{\texttt{common.py}, \texttt{merge.py} i ostali}
\begin{itemize}[nosep]
  \item \fajl{common.py} --- zajednički alati skrapera: \kod{Headline} struktura,
    HTTP \kod{HEADERS}, \kod{save\_csv}, konstanta putanje \kod{DATA\_RAW}.
  \item \fajl{merge.py} --- spaja više povlačenja i radi \textbf{deduplikaciju}:
    egzaktnu (identičan tekst) i \emph{približnu} (near-duplicate, biblioteka
    \kod{rapidfuzz}) da bi uklonio gotovo iste naslove.
  \item \fajl{rss\_scraper.py} --- rezervni izvor preko RSS-a.
  \item \fajl{mozzart\_bulk.py} --- skraper za Mozzart; \emph{napušten} (rad samo sa
    SportKlub-om).
\end{itemize}

% ===========================================================================
\section{Faza 2 --- Anotacija (\texttt{src/annotation/})}
Cilj faze: od sirovih naslova napraviti \textbf{balansiran, označen} skup i izmeriti
koliko se dva anotatora slažu (\textbf{Cohen's kappa}).

\subsection{\texttt{make\_splits.py} --- podela posla}
Deli skup \emph{tačno na pola} (Filip / Danilo, deterministički \kod{seed=42}). Da bi
se mogla izračunati saglasnost, izdvaja \textbf{kalibracioni skup}: prvih 110 naslova
svake polovine anotiraju \emph{oba} člana (unakrsno), pa za 220 naslova postoje dve
nezavisne oznake.

\begin{code}
rnd.shuffle(rows)                  # deterministički (seed 42)
A, B = rows[:pola], rows[pola:]    # Filipova / Danilova polovina
cal_A, cal_B = A[:110], B[:110]    # kalibracioni (prvih 110 svake) -> anotira i drugi
\end{code}

\textbf{Izlaz:} \kod{data/annotated/\{filip,danilo\}\_glavna.csv} (sa praznom kolonom
\kod{labela}) i \kod{data/calibration/*\_unakrsna.csv} (kalibracioni za drugog člana).

\subsection{\texttt{iaa.py} --- Cohen's kappa}
Računa međuanotatorsku saglasnost. \textbf{Cohen's kappa} meri slaganje \emph{iznad}
slučajnog: ako se dva anotatora slažu samo zato što obojica često biraju istu klasu,
kappa to „kažnjava''. Formula:
\[
\kappa = \frac{P_o - P_e}{1 - P_e},
\]
gde je $P_o$ stvarni udeo slaganja, a $P_e$ \emph{očekivano} slaganje slučajem.
Implementacija je svega nekoliko linija:

\begin{code}
def cohen_kappa(a, b):
    n = len(a)
    po = sum(1 for x, y in zip(a, b) if x == y) / n          # stvarno slaganje
    pe = sum((a.count(c)/n) * (b.count(c)/n)                  # očekivano slučajem
             for c in set(a) | set(b))
    return (po - pe) / (1 - pe), po
\end{code}

\textbf{Primer (naši brojevi).} Od 220 kalibracionih, slažu se na 182 $\Rightarrow$
$P_o = 0{,}827$. Filip je 93 naslova označio kao klikbejt (127 kao regularan), Danilo
81 (139 regularan). Očekivano slaganje:
\[
P_e = \tfrac{127}{220}\!\cdot\!\tfrac{139}{220} + \tfrac{93}{220}\!\cdot\!\tfrac{81}{220}
= 0{,}520.
\]
\[
\kappa = \frac{0{,}827 - 0{,}520}{1 - 0{,}520} = \mathbf{0{,}640}.
\]
Skripta ispiše i interpretaciju po Landis--Kohovoj skali (0,61--0,80 = „dobro'') i
listu svih neslaganja --- ulaz za analizu spornih slučajeva u izveštaju.

\begin{teorija}[ --- zašto kappa, a ne prosto slaganje]
Prosto procentualno slaganje ($P_o$) precenjuje kvalitet kada je zadatak
\emph{neuravnotežen} ili anotatori imaju naviku da biraju istu klasu. Ako bi npr.\ oba
anotatora nasumično pogađala, ali 80\,\% naslova nazivala regularnim, slučajno bi se
složili u $\approx 0{,}68$ slučajeva --- a to nije nikakvo „znanje''. \textbf{Cohen's
kappa} oduzima to očekivano slaganje slučajem ($P_e$) i normira ostatak:
$\kappa=1$ je savršeno slaganje, $\kappa=0$ je „kao slučaj'', $\kappa<0$ gore od
slučaja. Zato je kappa standard za međuanotatorsku saglasnost: meri slaganje
\emph{koje se ne može objasniti pukom raspodelom klasa}. Za binarne oznake u
ozbiljnim primenama poželjno je $\kappa\ge 0{,}8$; naših $0{,}640$ je „dobro'' i
posledica suštinske subjektivnosti pojma klikbejt (detaljno u izveštaju).
\end{teorija}

\subsection{\texttt{build\_dataset.py} --- balansiranje}
Spaja dve označene polovine i pravi \textbf{balansiran} finalni skup: deterministički
izmeša i uzme do 1100 naslova svake klase (1100 klikbejt / 1100 regularan):

\begin{code}
rnd = random.Random(42)
rnd.shuffle(sel_kb); rnd.shuffle(sel_ne)   # po klasi
bal = sel_kb[:1100] + sel_ne[:1100]
rnd.shuffle(bal)                            # promešaj klase
\end{code}

\textbf{Izlaz:} \kod{data/annotated/dataset.tsv} (format \kod{naslov\textbackslash t
labela} --- ulaz za Fazu 3) i nebalansirani \kod{dataset\_full.csv} (zadržava
\kod{id}/\kod{podtip}). Balansiranje uklanja pristrasnost klasifikatora ka većinskoj
klasi i pojednostavljuje tumačenje metrika.

\subsection{\texttt{reconcile\_iaa.py}, \texttt{stats.py}}
\begin{itemize}[nosep]
  \item \fajl{reconcile\_iaa.py} --- pomaže pri \emph{adjudikaciji} (usaglašavanju)
    spornih naslova posle IAA.
  \item \fajl{stats.py} --- deskriptivna statistika finalnog skupa (raspodela klasa,
    prosečne dužine) $\to$ \kod{results/faza2\_statistika.txt}.
\end{itemize}

% ===========================================================================
\section{Pretprocesiranje teksta (\texttt{src/preprocessing/serbian.py})}
Centralni modul za \textbf{tokenizaciju, stemovanje i lematizaciju}; koristi ga Faza
3a. Detaljno objašnjenje tehnika sa primerima je u zasebnom dokumentu
\texttt{docs/Obrada-teksta.md}; ovde sažeto i tehnički.

\textbf{Lanac:} \kod{basic\_normalize} (NFC + ćirilica$\to$latinica + lowercase)
$\to$ \kod{tokenize} $\to$ opciono \kod{stem}/\kod{lemma}.

\textbf{Tokenizacija} regularnim izrazom --- token je niz slova (sa srpskim
dijakriticima) i cifara; interpunkcija je razdvajač:
\begin{code}
# token = cifre + latinica + srpski dijakritici (zćčđš) + ćirilica (opseg A-S)
_TOKEN_RE = re.compile(r"[0-9a-zžćčđš<ćirilica>]+")
"Partizan pobedio 2:1!" -> ["partizan", "pobedio", "2", "1"]
\end{code}

\textbf{Stemovanje} --- heuristički, bez zavisnosti: za token se pohlepno (najduži
prvo) uklanja najduži poznati srpski nastavak, uz zaštitu da osnova ne padne ispod 3
slova:
\begin{code}
for suf in _SUFFIXES:                       # sortirani od najdužeg
    if token.endswith(suf) and len(token)-len(suf) >= 3:
        return token[:-len(suf)]
"pobedom","pobede","pobedi" -> "pobed"      # ista osnova
\end{code}

\textbf{Lematizacija} --- prava, preko biblioteke \kod{classla} (model za srpski,
procesori \kod{tokenize,pos,lemma}). Za razliku od stemera, vraća rečničku lemu i
hvata nepravilne oblike (\kod{najbolji}$\to$\kod{dobar}). Lazy-inicijalizacija i
keširanje modela, batch obrada radi brzine.

\textbf{Bez curenja informacija:} pošto su stem/lemma deterministički po tokenu,
pretprocesiranje se radi \emph{jednom} nad celim skupom; vokabular i IDF se uče tek
unutar svakog fold-a (vidi sledeću sekciju).

\begin{teorija}[ --- zašto uopšte normalizovati reči]
Klasični modeli vide tekst kao \textbf{vreću reči} (engl.\ \emph{bag of words}): svaka
različita reč je zasebno obeležje (kolona). Srpski je \emph{morfološki bogat} --- ista
reč ima desetine oblika (\kod{pobeda, pobede, pobedom, pobedi}\dots). Bez normalizacije
model svaki oblik uči zasebno, pa se signal \emph{raspršuje} po mnogo retkih kolona i
modelu brzo ponestane primera po obeležju (problem \emph{retkosti}/sparsnosti).
Stemovanje i lematizacija spajaju te oblike u jednu osnovu $\Rightarrow$ manje
dimenzija, gušće kolone, bolja generalizacija. \textbf{Stem vs lema:} stem je brza
heuristika (može pogrešiti), lema je morfološki tačna ali traži model --- zato obe
poredimo kao varijante. Lowercasing i transliteracija ćirilice su ista ideja:
spajanje površinski različitih, a značenjski istih tokena.
\end{teorija}

% ===========================================================================
\section{Faza 3a --- Osnovni modeli (\texttt{src/baseline/})}

\subsection{\texttt{run\_baseline.py} --- grid eksperimenata}
Poredi \textbf{logističku regresiju} i \textbf{multinomijalni naivni Bajes} preko
\kod{scikit-learn}, nad svim kombinacijama: normalizacija
(\kod{none}/\kod{stem}/\kod{lemma}) $\times$ ponderisanje
(\kod{bow}/\kod{tf}/\kod{idf}/\kod{tfidf}) $\times$ n-gram (1--1 / 1--2).

\textbf{Šeme ponderisanja} se mapiraju na parametre sklearn vektorizatora:
\begin{code}
bow   : CountVectorizer()                              # puko brojanje
tf    : TfidfVectorizer(use_idf=False, norm="l2")      # term-frekvencija
idf   : TfidfVectorizer(binary=True, use_idf=True)     # prisustvo × IDF
tfidf : TfidfVectorizer(use_idf=True, sublinear_tf=True, norm="l2")  # pun TF-IDF
# zajedničko: token_pattern=r"\S+"  (tekst je već tokenizovan), min_df=2
\end{code}
\begin{teorija}[ --- TF, IDF i TF-IDF ponderisanje]
Pošto reči nisu jednako informativne, broj pojavljivanja se \emph{pondériše}:
\begin{itemize}[nosep]
  \item \textbf{TF} (\emph{term frequency}) --- koliko se reč javlja u naslovu. Česta
    reč u dokumentu je važnija za taj dokument. (\kod{norm="l2"} skalira vektor na
    jediničnu dužinu da duži naslovi ne dominiraju.)
  \item \textbf{IDF} (\emph{inverse document frequency}) --- $\log(N/df)$: reč koja se
    javlja u \emph{svakom} naslovu (npr.\ „je'', „u'') nosi malo informacije i dobija
    malu težinu; retka, specifična reč dobija veliku.
  \item \textbf{TF-IDF} = TF $\times$ IDF --- spaja oba: visoko vrednuje reč koja je
    \emph{česta u ovom} naslovu, a \emph{retka u korpusu} (tj.\ diskriminativna).
    \kod{sublinear\_tf} primeni $1+\log(\mathrm{TF})$ jer 10 pojava nije 10$\times$
    važnije od jedne.
\end{itemize}
Zašto baš ove varijante poredimo: na \emph{kratkim} naslovima reči se retko ponavljaju,
pa TF$\approx$prisustvo i IDF nosi glavninu razlike --- empirijski proveravamo koliko
to vredi. \kod{min\_df=2} izbacuje reči iz samo jednog naslova (verovatno šum/tipfeler,
i sprečavaju preprilagođavanje). \textbf{N-grami} (1--2) uz pojedinačne reči hvataju i
kratke fraze („velika vest'') koje kao signal vrede više od zbira reči.
\end{teorija}
Bitno: \kod{lowercase=False} i \kod{token\_pattern=\textbackslash S+} jer je tekst
\emph{već} tokenizovan u \kod{serbian.py} --- vektorizator samo deli po razmaku, ne
tokenizuje ponovo.

\textbf{Pošteno ocenjivanje --- ugnežđena unakrsna validacija.} Da se izbegne curenje
informacija, koristi se \emph{nested CV}: spoljna 10-struka petlja meri performansu, a
unutrašnja bira hiperparametre (\kod{GridSearchCV}). Vektorizator je deo
\kod{Pipeline}-a, pa se vokabular/IDF uče \emph{samo} na trening delu svakog fold-a:

\begin{code}
outer = StratifiedKFold(10, shuffle=True, random_state=42)
for tr, te in outer.split(X, y):
    pipe = Pipeline([("vec", make_vectorizer(...)), ("clf", est)])
    gs = GridSearchCV(pipe, grid, scoring="f1", cv=inner)   # bira C / alpha
    gs.fit(X[tr], y[tr])                # uči SAMO na trening foldu
    y_pred = gs.best_estimator_.predict(X[te])
\end{code}

\begin{teorija}[ --- zašto baš \emph{ugnežđena} i \emph{stratifikovana} CV]
\textbf{Unakrsna validacija (CV)} deli skup na $k$ delova (fold-ova): trenira na
$k\!-\!1$, testira na 1, i tako $k$ puta --- pa usrednji. Time se performansa meri na
\emph{svim} podacima, a ne na jednom srećnom/nesrećnom razdvajanju.
\textbf{Stratifikovana} znači da svaki fold čuva odnos klasa (50/50) --- inače bi neki
fold mogao slučajno dobiti previše jedne klase i iskriviti metriku.
\textbf{Zašto \emph{ugnežđena} (nested):} ako biramo hiperparametre i \emph{istovremeno}
prijavljujemo rezultat na istom test skupu, „provirili'' smo u test $\Rightarrow$
optimistično pristrasna ocena. Zato unutrašnja petlja bira \kod{C}/\kod{alpha} samo na
trening delu, a spoljna petlja (koja test fold nikad nije videla pri izboru) daje
\emph{poštenu} ocenu. \kod{random\_state=42} čini podele ponovljivim.
\end{teorija}

\textbf{Hiperparametri --- čemu služe i zašto ti opsezi.} Hiperparametri se ne uče iz
podataka nego se zadaju i biraju validacijom. Mi tražimo po log-skali (geometrijski
korak), jer njihov efekat je multiplikativan, a ne linearan:
\begin{code}
LR  (logistička regresija):  C     in {0.01, 0.1, 1, 10, 100}
MNB (naivni Bajes):          alpha in {0.01, 0.1, 0.5, 1, 2}
\end{code}

\begin{teorija}[ --- C (regularizacija) i alpha (izglađivanje)]
\textbf{\kod{C} kod logističke regresije} kontroliše \emph{regularizaciju} --- kaznu za
prevelike težine. Model minimizuje (greška) $+\frac{1}{C}\lVert w\rVert$. \emph{Malo
C} $\Rightarrow$ jaka kazna $\Rightarrow$ jednostavniji model, manje
preprilagođavanja (\emph{overfitting}), ali rizik od \emph{podprilagođavanja}. \emph{Veliko
C} $\Rightarrow$ slaba kazna $\Rightarrow$ model se pripije uz trening podatke (rizik
overfittinga). Pošto se prava vrednost ne zna unapred, probamo opseg
$0{,}01\!-\!100$ i pustimo validaciju da izabere. Kod kratkih, retkih tekstova
regularizacija je važna jer obeležja ima mnogo, a primera malo.

\medskip
\textbf{\kod{alpha} kod naivnog Bajesa} je \emph{Laplace/aditivno izglađivanje}.
Naivni Bajes množi verovatnoće reči; ako se neka reč u trening skupu nikad nije pojavila
uz klasu, njena verovatnoća je 0 i \emph{poništi} ceo proizvod. \kod{alpha} dodaje malu
„pseudo-frekvenciju'' svakoj reči da nijedna ne bude nula. \emph{Veće alpha} =
glađe/opreznije procene (više se oslanja na prior), \emph{manje alpha} = vernije
trening podacima. Opet biramo validacijom u $0{,}01\!-\!2$.

\medskip
\textbf{Zašto log-skala:} skok sa $C{=}1$ na $2$ je beznačajan, ali sa $1$ na $100$
menja ponašanje modela --- zato koraci idu $\times10$, a ne $+1$ (efikasnije pokriva
prostor sa malo tačaka).
\end{teorija}

\begin{teorija}[ --- dva modela: zašto upravo NB i LR]
\textbf{Multinomijalni naivni Bajes} je verovatnosni klasifikator: pretpostavlja da su
reči \emph{nezavisne} u datoj klasi (otud „naivni'') i računa koja klasa čini naslov
verovatnijim. Brz, jak na malom tekstu, klasičan baseline za klasifikaciju teksta.
\textbf{Logistička regresija} je linearni model koji uči težinu svake reči i kombinuje
ih u verovatnoću klikbejta; ne pretpostavlja nezavisnost i daje dobro kalibrisane
verovatnoće. Dva komplementarna pristupa daju pošteniju „donju granicu'' koju
transformeri treba da nadmaše.
\end{teorija}

\subsection{\texttt{common.py} --- učitavanje i metrike}
\kod{load\_dataset} čita \kod{dataset.tsv}; \kod{compute\_metrics} računa iste metrike
za sve modele (radi poređenja): tačnost, preciznost/odziv/\textbf{F1 za klikbejt
klasu}, makro-F1 i \textbf{ROC-AUC} (kada model daje verovatnoće).

\begin{teorija}[ --- preciznost, odziv, F1, ROC-AUC]
\textbf{Preciznost} = od svih koje je model proglasio klikbejtom, koliko ih je stvarno
klikbejt (koliko „lažnih uzbuna''). \textbf{Odziv} (\emph{recall}) = od svih stvarnih
klikbejtova, koliko ih je model uhvatio (koliko je „propustio''). Njih dvoje su u
tenziji --- agresivniji model ima veći odziv ali manju preciznost. \textbf{F1} je
njihova harmonijska sredina, jedan broj koji nagrađuje balans (kažnjava ako je bilo
koje od dvoje nisko). \textbf{Zašto F1 baš na klikbejt klasi:} klikbejt (oznaka 1) je
klasa koja nas zanima da detektujemo; sama tačnost može biti varljiva, a F1 na ciljnoj
klasi direktno meri kvalitet detekcije. \textbf{Makro-F1} usrednji F1 obe klase
(tretira ih ravnopravno). \textbf{ROC-AUC} meri koliko dobro model \emph{rangira}
naslove po verovatnoći (nezavisno od praga 0,5): 1,0 savršeno, 0,5 nasumično ---
zato ga imaju samo modeli koji daju verovatnoću (NB, LR, BERT), a ne dekoder koji
vraća tvrdu 0/1 odluku.
\end{teorija}

\textbf{Izlaz:} \kod{results/baseline\_results.csv}. Najbolje: \emph{naivni Bajes +
stem + TF-IDF (1--2)}, F1$\approx$0,646.

% ===========================================================================
\section{Faza 3b --- Enkoderski LLM (\texttt{src/transformers/})}

\subsection{\texttt{finetune.py} --- fine-tuning BERTić i mBERT}
Fino podešava dva enkoderska transformera preko HuggingFace \kod{Trainer}-a:
\textbf{BERTić} (\kod{classla/bcms-bertic}, monolingvalni za BCMS) i \textbf{mBERT}
(\kod{bert-base-multilingual-cased}, višejezični). Protokol: 10-struka CV, variranje
broja epoha (2/3/4). Traži GPU --- pokreće se na Colab-u.

\begin{teorija}[ --- šta je BERT, enkoder vs dekoder, fine-tuning]
\textbf{Transformer} je neuronska mreža koja čita ceo tekst odjednom i preko mehanizma
\emph{pažnje} (\emph{attention}) za svaku reč gleda \emph{kontekst} (ostale reči). Otud
„kontekstualni'' modeli: ista reč dobija različitu reprezentaciju zavisno od okoline ---
za razliku od vreće reči koja je slepa za kontekst. \textbf{Enkoder} (BERT) je naučen da
\emph{razume} tekst (predviđa zamaskirane reči gledajući levo \emph{i} desno) i idealan
je za klasifikaciju. \textbf{Dekoder} (GPT, Claude) je naučen da \emph{generiše} tekst
(predviđa sledeću reč) --- koristimo ga u Fazi 3c.
\textbf{Fine-tuning} = uzmemo model već \emph{predtreniran} na ogromnom korpusu i
nastavimo obuku kratko na našem malom označenom skupu, da se prilagodi baš zadatku
klikbejta. \textbf{Mono vs multi:} BERTić je treniran samo na srpskom/srodnim jezicima,
pa bolje hvata nijanse od višejezičnog mBERT-a koji deli kapacitet na 100+ jezika ---
zato BERTić pobeđuje.
\end{teorija}

\textbf{Tehnički zanimljiv detalj --- sudar imena.} Folder se zove \kod{transformers},
isto kao HuggingFace biblioteka. Da \kod{import transformers} ne učita \emph{lokalni}
folder umesto biblioteke, kod na startu čisti \kod{sys.path} i keš modula:
\begin{code}
sys.path[:] = [p for p in sys.path if p not in ("", _SELF, _SRC)]
sys.path.append(_SRC)   # 'src' ide na KRAJ (da se 'baseline' i dalje nađe)
# + izbaci eventualno već uvezen lokalni 'transformers' iz sys.modules
\end{code}

\textbf{Tokenizacija za transformer} je drugačija od baseline-a: koristi se
\emph{subword} tokenizator samog modela (\kod{AutoTokenizer}), sa odsecanjem na
\kod{max\_len}. Za svaki fold pravi se \emph{svež} model i trenira tačno $n$ epoha:
\begin{code}
tokenizer = AutoTokenizer.from_pretrained(hf_id)
model = AutoModelForSequenceClassification.from_pretrained(hf_id, num_labels=2)
\end{code}

\begin{teorija}[ --- subword tokenizacija i hiperparametri obuke]
\textbf{Subword tokenizacija} deli reč na česte pod-delove (npr.\ \kod{pobedom}$\to$
\kod{pobed}+\kod{\#\#om}). Time model nema „nepoznatih reči'': svaku, i retku, sastavi
od delova --- posebno korisno za morfološki bogat srpski. Zato ovde \emph{ne} radimo
stemovanje: model sam uči podreči. Ključni hiperparametri obuke:
\begin{itemize}[nosep]
  \item \textbf{Broj epoha} (2/3/4) --- koliko puta model prođe ceo trening skup.
    Premalo $\Rightarrow$ ne nauči (podprilagođavanje); previše $\Rightarrow$ zapamti
    trening (overfitting). Zato \emph{variramo} i biramo: 3--4 je optimum (kriva učenja
    se izravna).
  \item \textbf{\kod{max\_len}} --- maksimalna dužina u tokenima; naslovi su kratki pa
    se gotovo ništa ne odseca, a računanje je brže.
  \item \textbf{\kod{batch\_size}} --- koliko primera ide kroz model pre ažuriranja
    težina; veći batch = stabilnije, ali traži više GPU memorije.
  \item \textbf{\kod{learning\_rate}} --- veličina koraka pri učenju; za fine-tuning je
    mali (npr.\ $2\!\cdot\!10^{-5}$) da se ne „pokvari'' predtrenirano znanje.
\end{itemize}
\textbf{Zašto svež model po fold-u:} da svaki fold bude nezavisan eksperiment, bez
prenosa naučenog iz prethodnog (inače curenje informacija).
\end{teorija}

\textbf{ROC-AUC} se dobija \emph{softmax}-om nad logitima (verovatnoća klikbejt klase):
\begin{code}
proba = softmax(logits)[:, 1]   # za roc_auc_score
\end{code}
\begin{teorija}[ --- logiti i softmax]
Model na izlazu daje \emph{logite} --- sirove brojeve za svaku klasu, koji nisu
verovatnoće. \textbf{Softmax} ih pretvara u verovatnoće (pozitivne, sabiraju se u 1).
Uzimamo verovatnoću klase 1 (klikbejt) da bismo mogli da računamo ROC-AUC, koji traži
rangiranje po verovatnoći, a ne samo tvrdu odluku.
\end{teorija}

\textbf{Izlaz:} \kod{results/encoder\_bertic\_results.csv},
\kod{encoder\_mbert\_results.csv}. Najbolje: BERTić, 3--4 epohe (F1$\approx$0,703,
ROC-AUC$\approx$0,788). \fajl{colab\_train.py} je verzija prilagođena Colab okruženju.

% ===========================================================================
\section{Faza 3c --- Dekoderski LLM (\texttt{src/decoder/})}

\subsection{\texttt{prompts.py} --- konstrukcija upita}
Definiše prompt varijante koje se porede: jezik (SR/EN) $\times$ format
(zero/few-shot) $\times$ sa/bez definicije klikbejta. \kod{build\_prompt} sklapa
\kod{(system, user\_template)}; ključno je da se od modela traži \textbf{samo cifra}:
\begin{code}
rule = "Odgovori ISKLJUČIVO jednom cifrom: 1 (klikbejt) ili 0 (regularan)."
user = f'{task}\n\n{definicija}{examples}Naslov: "{{naslov}}"\nOdgovor:'
\end{code}
Few-shot varijante ubacuju 6 rešenih primera (\kod{FEWSHOT}); ti naslovi se po potrebi
izuzimaju iz evaluacije (holdout) da ne bi „procurili''.

\begin{teorija}[ --- zero-shot vs few-shot, uloga definicije]
\textbf{Zero-shot} = modelu damo samo zadatak, bez ijednog rešenog primera --- oslanja
se isključivo na znanje iz predtreniranja. \textbf{Few-shot} = u prompt ubacimo
nekoliko rešenih primera (ovde 6) iz kojih model „na licu mesta'' nauči obrazac
(\emph{in-context learning}), bez menjanja težina. Zato few-shot primeri \emph{ne} smeju
biti u test skupu --- inače merimo memorisanje, ne generalizaciju (otud \emph{holdout}).
\textbf{Definicija klikbejta} u promptu je treći faktor: usmerava model na \emph{naše}
shvatanje pojma. Sve tri ose (jezik, broj primera, definicija) variramo da vidimo šta
najviše pomaže --- to je „prompt inženjering'' kao zamena za obuku.
\end{teorija}

\subsection{\texttt{claude\_eval.py} --- evaluacija nad celim skupom}
Šalje svaki naslov Anthropic Claude-u (Haiku), zero-shot, i parsira 0/1 odgovor.
\kod{max\_tokens=16} jer je odgovor jedna cifra (jeftino i brzo):
\begin{code}
resp = client.messages.create(model=model, max_tokens=16,
        system=system, messages=[{"role": "user", "content": user_prompt}])
\end{code}
\begin{teorija}[ --- zašto \kod{max\_tokens=16} i bez \emph{temperature}]
\kod{max\_tokens} ograničava dužinu odgovora; pošto tražimo samo cifru 0/1, 16 je
sasvim dovoljno i \emph{jeftinije} (plaća se po tokenu). \emph{Temperature} kontroliše
nasumičnost generisanja --- za klasifikaciju želimo \emph{determinističan}, ponovljiv
odgovor, pa je ne dižemo (na nekim modelima bi je slanje čak odbilo). Sistemski prompt
izričito traži samo cifru da bi parsiranje (\kod{parse\_label}) bilo pouzdano.
\end{teorija}

\textbf{Tri inženjerska detalja} koja čine evaluaciju pouzdanom i jeftinom:
\begin{enumerate}[nosep]
  \item \textbf{Keširanje} (\kod{results/decoder\_cache/*.jsonl}) --- svaki rezultat se
    upiše; ponovno pokretanje preskače već klasifikovane naslove (prekid/nastavak bez
    dvostrukog plaćanja).
  \item \textbf{Retry sa eksponencijalnim odlaganjem} --- na privremenu grešku/rate
    limit pokušava do 5 puta (\kod{2**attempt}).
  \item \textbf{Ograničavanje učestalosti} (\kod{min\_interval}) --- razmak između
    poziva, da se ne probije API limit.
\end{enumerate}

\textbf{Izlaz:} \kod{results/decoder\_results.csv} + log. Najviši F1 na klikbejt klasi
u celom radu ($\approx$0,715), ali samo tvrde 0/1 odluke (bez ROC-AUC).
\fajl{gemini\_eval.py} i \fajl{chatgpt\_eval.py} su alternativni dekoderi (dele keš i
parsiranje).

% ===========================================================================
\section{Faza 4 --- Izveštaj (\texttt{src/report/} + \texttt{report/})}

\subsection{\texttt{make\_tables.py} --- CSV $\to$ LaTeX tabele}
Čita \kod{results/*.csv} i piše \kod{booktabs} tabele spremne za \kod{\textbackslash
input} u izveštaj (\kod{report/tables/}). Bira najbolju varijantu po tipu modela za
zbirnu tabelu poređenja. Otporno na nedostajuće ulaze --- ako CSV ne postoji, upiše
placeholder da se izveštaj svejedno kompajlira.

\subsection{\texttt{make\_figures.py} --- grafici}
\kod{matplotlib} grafici: raspodela klasa (bar), histogram dužine naslova po klasi,
kriva učenja (F1/AUC po epohi) i poređenje modela. Svaki se snima kao PNG (150 dpi) u
\kod{report/figures/}.

\subsection{\texttt{izvestaj.tex} --- glavni dokument}
LaTeX izveštaj (pdfLaTeX / tectonic); \kod{\textbackslash input}-uje generisane tabele
i uključuje grafike. Komajlira se u \kod{report/izvestaj.pdf}. Bibliografija je u
\kod{references.bib}.

% ===========================================================================
\section{Kako sve teče zajedno (reprodukcija)}
Ceo lanac, ukratko (pune komande sa očekivanim izlazom su u
\texttt{docs/Komande.md}):
\begin{enumerate}[nosep]
  \item \kod{sportklub\_api.py} $\to$ sirovi naslovi.
  \item \kod{make\_splits.py} $\to$ (ručna anotacija) $\to$ \kod{iaa.py} (kappa) $\to$
    \kod{build\_dataset.py} $\to$ \kod{dataset.tsv}.
  \item \kod{run\_baseline.py}, \kod{finetune.py}, \kod{claude\_eval.py} $\to$
    \kod{results/*.csv}.
  \item \kod{make\_tables.py} + \kod{make\_figures.py} $\to$ \kod{tectonic
    izvestaj.tex} $\to$ \textbf{PDF}.
\end{enumerate}

\vspace{2mm}\hrule\vspace{2mm}
\noindent\textit{Prateći dokumenti:} \texttt{docs/Mapa-koda.md} (kratka mapa),
\texttt{docs/Obrada-teksta.md} (tokenizacija/stemovanje/lematizacija detaljno),
\texttt{docs/Komande.md} (komande + očekivani izlaz).

\end{document}
